\section{Load fragments when a consumer needs them}
\label{sec:pg-storage-reading}

Reading a parent table and reading sequence symbols are separate steps:

\begin{enumerate}
  \item Read visible stripe and chunk metadata under the scan snapshot.
  \item Select the required columns and reject chunk groups contradicted
    by supported scalar bounds.
  \item Read the selected presence and stored value buffers, decode a
    group, and apply row masks and remaining conditions.
  \item Bind the surviving row's locator to its source, validating the
    source identity, epoch, and fragment geometry under the scan snapshot.
  \item When symbols are requested, fetch the fragment rows covering
    that interval and check their packed representation.
\end{enumerate}

Chernishev et al. study \emph{late materialization}, which uses row
positions to postpone tuple reconstruction~\autocite{Chernishev2022ACS}.
In BioSeqDB's lazy read path, a tuple carries a sequence locator, and
fragments are fetched when an operation needs their symbols. Sequence
contents can therefore remain unread after PostgreSQL has formed a tuple
containing the locator.

The column reader loads stored buffers for all selected groups of a stripe
before decoding groups in turn. Thus avoiding a whole-sequence allocation
does not mean that the complete scan retains only one compressed chunk.
Likewise, a fragment request can cause its containing column chunk to be
read and decoded. The logical request and the physical I/O unit differ.

The table reader can return ordinary tuple slots or column batches containing
projected values, NULL masks, a mask of surviving rows, and internal row
numbers. Supported scalar conditions can update a column batch's mask.
Surviving rows can then be extracted into ordinary tuple slots. Alignment
separately collects its input rows into the batches described in
Section~\ref{sec:customscan-batch-state-machine}. These are different work
units. Paths that pass source locators through or evaluate sequence
expressions retain scalar fetch where the vector conditions are disabled.
Batching itself does not imply a hardware SIMD instruction for every
expression.

How a scanner copies a fragment into its own slots, how long those slots
live, and how that copy sits next to the column chunk cache is described
in Section~\ref{sec:lazy-sequence-facade}.

\subsection{Decoded column chunks have a different cache}
\label{sec:pg-storage-caches}

The optional column decompression cache stores decoded compressed chunks
inside one PostgreSQL backend. Its key contains relation, stripe, chunk
group, and column identifiers. It is disabled by default; the configured
capacity when enabled defaults to 200~MiB.

The cache counts reads when selecting entries to remove. It protects one
current chunk group per relation from removal. A cache hit returns a borrowed
buffer, which must not outlive its cache entry. Uncompressed chunks do not
need this decompression cache. Reuse can save repeated decoding during a scan
or rescan before scan cleanup clears the cache. It does not share a decoded
buffer between sessions.

\begin{table}[htbp]
\centering
\small
\caption{Stored bytes and reusable memory have different owners.}
\label{tab:storage-cache-scopes}
\begin{tabularx}{\linewidth}{@{}p{.23\linewidth}XX@{}}
\toprule
\textbf{Object} & \textbf{Retained data} & \textbf{Scope}\\
\midrule
PostgreSQL buffer & Relation page bytes & Shared database buffer manager;
short pins protect each page operation\\
Fragment slot & Copy of one packed fragment & Sequence scanner and its
read-session owner\\
Column cache & Decoded column value buffer & One backend, subject to its
configured byte limit\\
Query-result cache & Completed SQL result rows & Separate extension shared
across backends\\
\bottomrule
\end{tabularx}
\end{table}

\section{Physical storage of shared query results}
\label{sec:pg-storage-request-coordination}

The separate \texttt{bioseq\_cache} extension stores completed SQL results,
not source fragments. Its producer and waiter lifecycle, exact matching,
and generic range-reuse contract are described in
Section~\ref{sec:bioseq-result-cache}. The physical store groups result rows
into columnar blocks with a schema description, NULL information, value
offsets, and checksums. Variable-length values are flattened before storage;
process-local pointers cannot become reusable result bytes.

Completed blocks can live in shared memory or cache-owned private files.
Before storing a block, the cache reserves capacity. For shared memory, it
checks the configured byte ceiling and the minimum amount of system memory
that must be left free. Before writing a file block, it also checks the file
capacity limit and the minimum amount of disk space that must be left free.
If neither tier can
accept another block, the original query continues without publishing a
reusable result. Active readers keep references that prevent their blocks
from being removed.

This file tier is derived data with its own cleanup rules. Cached results
are cleared across a server restart. It should not be confused with the
PostgreSQL relations holding authoritative sequence data or with persistent
sequence indexes.

\section{From stable sources to persistent indexes}
\label{sec:pg-storage-summary}

The storage layer lets two queries use the same table in different ways.
A metadata query reads selected column chunks and can avoid sequence access
entirely. An alignment query keeps a validated locator and asks for packed
fragments under the same snapshot as its input row. Updating or deleting the
row changes what newer snapshots may read while preserving the old source
for an existing reader.

The next chapter builds reusable lookup structures on these sources. It
uses the same distinction between a location, a source version, and its
lifetime. An index can still occupy bytes after its source is invalidated;
being present in storage is not sufficient permission to use it.
